\section{The synthetic contact binaries}
\label{app:synth}

This appendix specifies the synthetic bodies used in
Section~\ref{sec:constrains:synthetic}, states the design decisions that make
them a controlled experiment rather than a demonstration, and records two
errors in the first version of that experiment.

\subsection{Why synthetic bodies}
\label{app:synth:why}

The background density, and with it the excess mass that
Eq.~\eqref{eq:excessidentity} ties to it, is the more invariant quantity on
Itokawa and the less invariant one on 67P and Arrokoth. Three real bodies establish that the
behaviour varies; they cannot establish what it varies with. A fourth real body
would not settle it either. We cannot choose in advance whether a given
asteroid carries a detectable anomaly, and no body's interior is known
independently, so agreement would establish only that something had occurred
twice.

What is needed is a family in which the strength of the signal can be set by
the experimenter while everything else is held fixed. That requires synthetic
bodies. The cost is that they are idealised, and Section~\ref{app:synth:limits}
states what that costs.

\subsection{Construction}
\label{app:synth:build}

Each body is the union of prolate lobes of revolution about $x$. The outer
half-width is
\begin{equation}
w(x) = \max_{i}\; A_i \sqrt{1 - \left(\frac{x - c_i}{a_i}\right)^{2}},
\label{eq:appsynthprofile}
\end{equation}
taken as zero where the radicand is negative, with $A_i$, $c_i$ and $a_i$ the
amplitude, centre and half-length of lobe $i$. The surface is meshed directly
by revolving the profile: $100$ stations in $x$, $48$ around, with apex
vertices closing both ends, giving about $9{,}500$ facets. Orientation is fixed
by requiring the total signed volume to be positive.

Using a union of analytic lobes rather than a boolean operation on two meshes
avoids the intersection curve entirely, so no body in the family contains a
seam, a sliver, or a non-manifold edge. It also makes each body a single
connected component, so the containment test of
Section~\ref{app:clip:verify} has nothing to find: the failure mode that
affects the Arrokoth shape model cannot arise here. The profile is trimmed at a half-width
of $10^{-6}$ of its maximum rather than at an absolute threshold; an absolute
one leaves a ring of vertices a few nanometres from the axis near each tip,
distinct by index but coincident in position, which is enough to make the mesh
fail a manifold test. Every body in the family, as built, has
$\chi = 2$ with no repeated or unmatched directed edges, and no clipped region
in the experiment was rejected by the manifold check.

\subsection{The family}
\label{app:synth:family}

The eight bodies fall into three groups, and the grouping is the design.

A near-spherical control stands in for a body with no interior structure. A
\emph{symmetric} family of equal lobes is symmetric about its own neck, so the
density that flattens the geopotential should be close to uniform and the
signal weak by construction, whatever the depth of the neck. An
\emph{asymmetric} family of unequal lobes has a geometric mass asymmetry for
the inversion to find. Within each family the neck is deepened by moving the
lobes apart while the enclosed volume varies by no more than seven per cent, so
that depth of neck and size of body are not confounded.

\begin{table}[htbp]
\centering
\small
\caption{The synthetic family. Lobes are given as $(A_i, c_i, a_i)$ in km.
Prominence is that measured by the detector of
Section~\ref{sec:methods:neck}; the neck ratio is the half-width at the saddle
divided by the largest half-width. All bodies were run at a bulk density of
$2000$~kg\,m$^{-3}$ and a rotation period of $12.13$~h, with seven planes
spanning a factor $1.900$ in region volume, and all have $\chi = 2$.}
\label{tab:appsynth}
\begin{tabular}{llrrr}
\toprule
Body & Lobes $(A, c, a)$ & $\Vtot$ (km$^3$) & Neck ratio & Prominence \\
\midrule
\texttt{sphere} & $(0.15, 0.00, 0.16)$ & $0.01503$ & $0.650$ & none \\
\addlinespace
\texttt{sym\_waist} & $(0.14, \mp0.11, 0.185)$ & $0.02705$ & $0.787$ & $0.170$ \\
\texttt{sym\_neck} & $(0.14, \mp0.15, 0.185)$ & $0.02951$ & $0.587$ & $0.360$ \\
\texttt{sym\_deep} & $(0.14, \mp0.18, 0.185)$ & $0.03025$ & $0.235$ & $0.627$ \\
\addlinespace
\texttt{asym\_waist} & $(0.15, -0.13, 0.20)$, $(0.12, +0.12, 0.17)$ & $0.02707$ & $0.653$ & $0.124$ \\
\texttt{asym\_neck} & $(0.15, -0.13, 0.20)$, $(0.12, +0.18, 0.17)$ & $0.02849$ & $0.485$ & $0.272$ \\
\texttt{asym\_deep} & $(0.15, -0.13, 0.20)$, $(0.12, +0.22, 0.17)$ & $0.02893$ & $0.289$ & $0.449$ \\
\texttt{asym\_extreme} & $(0.15, -0.13, 0.20)$, $(0.10, +0.20, 0.14)$ & $0.02461$ & $0.196$ & $0.362$ \\
\bottomrule
\end{tabular}
\end{table}

The bodies are on the scale of Itokawa, a few hundred metres, and share its
rotation period. The near-spherical control has no saddle and its plane was
imposed, as Bennu's was.

\subsection{Selecting the planes}
\label{app:synth:planes}

Planes are chosen by bisection on the fraction of volume they leave above them,
not by position. For a target fraction $f$ the plane is found by bisecting on
$x$ until $\lvert V(x > \xsplit)/\Vtot - f \rvert$ is below tolerance; thirty
iterations give agreement to four decimal places in the fraction, at a cost of
under a second per plane.

Seven targets are placed geometrically about the fraction at the detected
saddle, spanning a factor
\begin{equation}
\frac{f_{\max}}{f_{\min}} = 1.900 ,
\label{eq:appspan}
\end{equation}
chosen to match the Itokawa sweep. This is the point on which the experiment
turns. A plane moved a fixed distance sweeps very different amounts of volume
on bodies of different shape---on Arrokoth, whose neck is long and thin,
$3.5$~km of travel changes the region volume by only five per cent against
ninety on Itokawa---so equal distances do not probe the same thing. Selecting
by volume fraction makes the sweeps comparable across bodies, and it is also
the variable against which the exponent of Eq.~\eqref{eq:slope} is fitted.

\subsection{What is not being tested}
\label{app:synth:not}

It is worth being explicit, since the natural expectation of a synthetic
experiment is different. We do not impose a known density on a synthetic body
and ask the method to recover it. That test would fail by construction and
would establish nothing.

Potential-dispersion minimization does not fit observed data. It assumes that
the true density is the one that renders the surface geopotential most uniform.
It assumes that the true density is the one that renders the surface
geopotential most uniform, and the objective is built from the shape and the
trial density alone: a field computed beforehand from an imposed interior is
never read. Planting a contrast in a body and leaving its shape unaltered
therefore changes nothing the minimization can see, and the recovered value is
whatever the shape alone would have given. What the experiment examines
is the structure of the objective function: which quantity it holds fixed as
the dividing plane moves, and how that depends on the strength of the response.

Every body in the family is therefore of uniform density. What varies is shape.
A consequence is that the experiment cannot establish that a strong response
implies a physical mass anomaly, only that in the strong-response regime the
quantity held fixed is the background density, and with it $\Dmass$, rather
than $\densR$. A genuine ground-truth test would have to build the body the
premise describes---a shape whose surface has relaxed onto an equipotential of
a planted interior---before inverting it.
Section~\ref{sec:discussion:rate} reports our attempt at that construction and
why its outcome is not presented.

\subsection{Results}
\label{app:synth:results}

\begin{table}[htbp]
\centering
\small
\caption{Results for the synthetic family. Half-ranges are taken over the seven
planes of each sweep, and standard errors on $s$ are those of the
least-squares fit. Where the density difference changes sign no exponent
exists and the fractional range of $\Dmass$ is not meaningful, its mean lying
inside its spread. $\bar{\ampfac}$ is the amplification factor of
Equation~\eqref{eq:excessamplify} averaged over the sweep: it is what
separates the $\denshead$ and $\Dmass$ columns, so those two are not
comparable across rows without it, and it is undefined on the symmetric
bodies, where $\Dmass$ passes through zero. The last column is the
compensation ratio, the half-range of $\Ddens$ divided by that of $\Dmass$,
which is the part of the stability ratio that is not fixed by algebra
(Section~\ref{sec:constrains:observation}). The Itokawa row is the eleven
planes below the grain-density ceiling, the set whose span of $1.907$ matches
the $1.900$ used here; its maximum improvement over that set is $22.4\%$
rather than the $16.9\%$ quoted at the detected neck.}
\label{tab:appsynthres}
\begin{tabular}{lrrrrrr}
\toprule
Body & $\improve_{\max}$ & $s$ & $\bar{\ampfac}$ & $\denshead$ half-range
 & $\Dmass$ half-range & $\Ddens/\Dmass$ \\
\midrule
\texttt{sphere}       & $1.24\%$  & none & $44$ & $3.10\%$  & sign change & --- \\
\texttt{sym\_waist}   & $3.68\%$  & none & ---  & $5.61\%$  & sign change & --- \\
\texttt{sym\_neck}    & $4.35\%$  & none & ---  & $6.22\%$  & sign change & --- \\
\texttt{sym\_deep}    & $4.47\%$  & none & ---  & $5.60\%$  & sign change & --- \\
\addlinespace
\texttt{asym\_waist}  & $16.67\%$ & $-2.020\pm0.109$ & $15.0$ & $9.28\%$  & $31.3\%$ & $1.9$ \\
\texttt{asym\_neck}   & $17.63\%$ & $-2.366\pm0.165$ & $16.2$ & $9.37\%$  & $39.2\%$ & $1.7$ \\
\texttt{asym\_deep}   & $19.50\%$ & $-2.069\pm0.094$ & $14.9$ & $8.81\%$  & $29.1\%$ & $2.0$ \\
\addlinespace
\texttt{asym\_extreme}& $43.71\%$ & $-1.217\pm0.023$ & $8.1$  & $12.78\%$ & $6.2\%$  & $6.1$ \\
\midrule
Itokawa (real)        & $22.4\%$  & $-1.141\pm0.018$ & $11.8$ & $13.20\%$ & $5.12\%$ & $7.1$ \\
\bottomrule
\end{tabular}
\end{table}

The design succeeded in separating the groups. Every symmetric body, however
deep its neck, produces a weak signal and a density difference that changes
sign across its sweep; the deepest, \texttt{sym\_deep}, has a prominence of
$0.627$---the highest in the family, a more sharply necked profile than
\texttt{asym\_extreme}---and still returns only
$4.47\%$. Depth of neck is not signal strength. The asymmetric bodies produce
signals three to ten times larger at comparable neck depths, because they have
a mass asymmetry to find.

The exponent varies monotonically with signal strength across the four bodies
that have one, and the Pearson correlation is $0.95$ ($p = 0.05$ two-tailed, $n = 4$), while
the Spearman rank correlation over the same four is only $0.40$, the three
clustered bodies not being in rank order. Only \texttt{asym\_extreme}
approaches $-1$, the value at which Eq.~\eqref{eq:slopeequiv} makes the
background density the quantity held fixed as the plane moves, and it is the only body in the family that reproduces Itokawa's
behaviour---on the exponent, on the density half-range, and on the excess-mass
half-range simultaneously.

\subsection{Two errors in the first version}
\label{app:synth:errors}

Both are recorded because either would have produced a plausible and wrong
conclusion.

\emph{The sweeps were not comparable.} The first version swept a fixed
$\pm0.05$~km on every body. Because the bodies differ in how sharply their
cross-section changes near the neck, this produced volume-fraction spans of
$1.13$ to $1.54$ against the $1.90$ of the Itokawa sweep. A body whose region
volume barely changes cannot show either quantity varying, so the comparison
measured the sweep design rather than the bodies. The bisection scheme of
Section~\ref{app:synth:planes} was introduced to correct this.

\emph{The diagnostic broke down where it mattered most.} The first version
compared fractional half-ranges directly. That statistic divides by the mean,
and on a body with no anomaly the mean of $\Dmass$ sits inside its own spread,
so the quotient is arbitrary. The symmetric bodies returned apparent ratios of
up to $1623$, which were read at first as an extraordinarily strong result in
the opposite direction. They were artefacts of dividing by nearly zero. The
exponent of Eq.~\eqref{eq:slope} was introduced because it does not have this
failure mode: where the density difference changes sign, the exponent simply
does not exist, and its absence is the correct answer rather than a large
number.

\subsection{Limitations}
\label{app:synth:limits}

Four bodies in the family carry a measurable exponent, and three of those
cluster between $16.7\%$ and $19.5\%$ in signal strength while the fourth sits
at $43.71\%$. The correlation therefore rests on a single high-leverage point,
which is what the Spearman value of $0.40$ records.
The direction of the effect is supported; its functional form is not
determined, and a threshold anywhere between roughly twenty and forty-four per
cent fits the data as well as a gradual trend does. Filling that gap would
require bodies designed to land in it, which the present family does not
attempt.

The numerical threshold does not transfer to real bodies. \texttt{asym\_extreme}
requires $43.71\%$ to reach the regime in which the background density is
pinned, while Itokawa reaches it at $22.4\%$ over the eleven planes below
the grain-density ceiling, or $26.4\%$ if the compositionally excluded
planes are counted. The synthetic surfaces are smooth surfaces of revolution; real
bodies carry roughness at every resolved scale, which Paper~IV found
contributes an error floor of order two percentage points to the dispersion. No
threshold value from this experiment should be quoted as universal, and none is
quoted in this paper.

The experiment uses the same forward model in both directions and therefore
examines the objective function, not the accuracy of the gravity calculation.
That accuracy is tested separately in
Section~\ref{sec:discussion:ellipsoid} against a homogeneous ellipsoid with a
closed-form surface gravity.

Finally, the family varies one geometric parameter---the separation of two
lobes of fixed shape. Real contact binaries differ in more ways than that, and
whether the relation between signal strength and exponent survives variation in
lobe elongation, in the ratio of lobe sizes, or in spin rate has not been
tested here.


